\section{Treatment: an identification-first causal pipeline}
\label{ann:treatment}

This baseline cohort has no randomization (the trial-arm field is a DSM-5 subtype), so
treatment \emph{selection}---who should get which drug---is a counterfactual question
observational treatment-as-usual cannot answer. The pipeline therefore \emph{bounds}
what the map licenses and \emph{defends} the prognostic finding, rather than claiming a
treatment effect.

\subsection{Exposures, overlap and weighting}

Treatment exposures, recovered from the per-cohort source dictionaries, were harmonized
to common drug-class indicators. For each treatment contrast we estimate a propensity
$e(v_i)=P(\text{treat}\mid v_i)$ on a design $v_i$ that includes severity, diagnosis,
demographics \emph{and} the nine coordinates~\citep{rosenbaum1983propensity}, gate on
common support (retain the region where both arms have mass), and reweight by stabilized
inverse-probability-of-treatment weights~\citep{austin2011iptw}
\begin{equation}
  w_i=\frac{P(\text{treat})}{e(v_i)}\ \text{(treated)},\qquad
  w_i=\frac{1-P(\text{treat})}{1-e(v_i)}\ \text{(control)} ,
\end{equation}
checking covariate balance by the post-weighting maximum standardized mean difference.
Lithium-in-BP balanced to a maximum SMD of $0.008$ (near-complete overlap $0.997$);
antipsychotic-in-BP balanced to $0.079$; clozapine-in-SZ retained residual imbalance
(SMD $0.335$), the signature of channelling.

\subsection{Doubly-robust moderation, the E-value and the minimum detectable effect}

The estimand of interest is \emph{moderation}: does a durable axis or archetype change
\emph{who} benefits. We fit a doubly-robust errors-in-variables outcome model with a
treatment-by-coordinate interaction (the EIV term as in \eqref{eqD:eiv}), reporting the
per-axis interaction credible interval and the held-out $\Delta\mathrm{ELPD}$ of adding
the interaction. Each average treatment effect carries an E-value for unmeasured
confounding~\citep{vanderweele2017evalue}---the minimum strength of association an
unmeasured confounder would need, with both treatment and outcome, to explain away the
estimate---and, crucially, a \emph{minimum detectable effect} (MDE): the smallest
interaction the design resolves at $80\%$ power. The MDE is what makes a null
\emph{bounded} rather than blind.
\begin{itemize}
  \item \textbf{Lithium-BP} (cleanest): average effect on functioning $-0.003$ (E-value
  $1.06$), with per-axis interactions well inside an MDE of $0.19$--$0.22$ SD---a
  \emph{well-identified, bounded null}. The design could have resolved a meaningful
  interaction and did not.
  \item \textbf{Antipsychotic-BP}: a confounded \emph{average} effect ($-0.231$, E-value
  $1.77$) but \emph{no reliable moderation}---the interaction $\Delta\mathrm{ELPD}$ band
  spans zero and the two map encodings disagree on which axis carries it (the durable
  representation points at metabolic, the archetype representation at the biological
  corner), textbook false-positive behaviour.
  \item \textbf{Clozapine-SZ}: channelled and underpowered (MDE $0.37$--$0.67$ SD), a
  non-decisive null rather than a clean one.
\end{itemize}
Average treatment effects were uniformly confounding-fragile (E-values $1.1$--$1.8$).

\subsection{The map is not a treatment proxy (defending the prognosis)}

The most damaging alternative to the prognostic finding is that the biological corner
simply received different drugs that drove the outcome. We refit the functioning
prognosis with and without the harmonized drug-class exposures, both unweighted and
under attrition weighting. The archetype carrier of the forecast survived
($\beta$ $0.164\!\to\!0.156$ unweighted, attenuation $4.7\%$;
$0.147\!\to\!0.141$ under attrition weighting, $3.9\%$; both credible intervals
excluding zero), so the functional forecast is not merely unmodelled treatment. (The
isolated durable trio did not survive unweighted, consistent with the prognosis annex;
its metabolic axis sharpened to exclude zero only under weighting.) The adjustment is
for baseline/lifetime (BP) or current (SZ/DR) drug-class exposure, not a marginal
structural model over time-varying treatment.

\subsection{Treatment-course atlas}

Even without selecting a drug, the baseline archetype \emph{describes} two-year
treatment course. Per corner (pooled BP$+$SZ, Wilson intervals): treatment-resistance
$43\%$ (biological) to $20\%$ (low-burden); CGI response $46\%$ to $59\%$; significant
side-effects $23\%$ to $11\%$. The corner added beyond baseline severity, substance
comorbidity and demographics (likelihood-ratio $p\le3\times10^{-3}$ for all three
endpoints), the gradient was within-cohort (composition share $\le6\%$,
corner-by-cohort interaction non-significant), and held-out $\Delta\mathrm{ELPD}$ was
$+20/+16/+10$. Individual discrimination was modest: under a permutation null on the
held-out $\Delta$AUC, response ($p=0.015$) and side-effects ($p=0.005$) cleared, while
treatment-resistance---the steepest gradient---was AUC-marginal ($p=0.185$), the same
group-versus-individual collapse seen for prognosis. The atlas is therefore a
monitoring signal (which phenotype to watch more closely), never a prescription.
